\begin{theorem}[Lee--Yang 特征曲线的椭圆化与 \(37\)-模性]\label{thm:xi-terminal-zm-leyang-elliptic-modularity-37}
设
\[
\Pi(\lambda,y):=\lambda^{4}-\lambda^{3}-(2y+1)\lambda^{2}+\lambda+y(y+1)\in\ZZ[\lambda,y],
\]
并令仿射曲线 \(\mathcal C\subset\mathbb A^2_{(\lambda,y)}\) 由 \(\Pi(\lambda,y)=0\) 给出。
令
\[
x:=\lambda,\qquad v:=y-x^{2}.
\]
其中 \(v\) 与前文记号 \(u:=y-\lambda^2\) 一致。
则 \(\mathcal C\) 与椭圆曲线
\[
E_0:\quad v^{2}+v=x^{3}-x
\]
在 \(\QQ\) 上同构；可取显式互逆同构
\[
\phi:E_0\to \mathcal C,\quad (x,v)\longmapsto(\lambda,y)=(x,v+x^{2}),
\qquad
\phi^{-1}:\mathcal C\to E_0,\quad (\lambda,y)\longmapsto (x,v)=(\lambda,y-\lambda^{2}).
\]
并且 \(\Delta(E_0)=37\)、\(j(E_0)=110592/37\)；在 Cremona/LMFDB 记号下 \(E_0\) 对应强 Weil 曲线 \(37\mathrm a1\)，并可作为
\(
X_0^+(37):=X_0(37)/\langle w_{37}\rangle
\)
的代数模型之一。\cite{LMFDBEllipticCurve37a1}
\end{theorem}
\begin{proof}
将 \(y=v+x^2\) 与 \(\lambda=x\) 代入得到恒等式
\[
\Pi(x,v+x^2)=v^2+v-x^3+x,
\]
从而 \(\Pi(\lambda,y)=0\) 与 \(v^2+v=x^3-x\) 等价；并且 \((x,v)\leftrightarrow(\lambda,y)\) 的互逆性由 \(y=v+x^2\) 在仿射坐标环中的可逆性直接给出，故 \(\phi\) 为 \(\QQ\)-同构。
不变量计算与外部标识见注记 \ref{rem:xi-terminal-elliptic-37a1-modular-anchor}。证毕。
\end{proof}
\par

\begin{proposition}[重量函数的分歧结构、切触多项式与临界点的群律约束]\label{prop:xi-terminal-zm-leyang-f-ramification}
令 \(E_0\) 为定理 \ref{thm:xi-terminal-zm-leyang-elliptic-modularity-37} 的椭圆曲线，\(O\) 为无穷远点。
设
\[
f:=v+x^{2}=\phi^{*}(y)\in\QQ(E_0),
\]
并以同一符号记由 \(f\) 诱导的有理映射 \(f:E_0\to\PP^{1}\)。
则：
\begin{enumerate}
  \item \(f\) 在 \(O\) 处具有唯一四阶极点，且 \((f)_\infty=4O\)。因此 \(\deg(f)=4\)，并在 \(O\) 上方全分歧（分歧指数 \(e_O=4\)）。
  \item \(f\) 的有限分歧点恰为 \(\dd f=0\) 的点；其 \(x\)-坐标满足
  \[
  \dd f=0\quad\Longleftrightarrow\quad (x^{2}-1)\bigl(16x^{3}-9x^{2}+1\bigr)=0.
  \]
  更具体地，有限分歧点由五点组成：两条有理点
  \[
  P_0=(1,-1),\qquad P_1=(-1,0),
  \]
  以及三点 \(P_2,P_3,P_4\)，其中 \(x(P_i)\) 为三次方程 \(16x^{3}-9x^{2}+1=0\) 的三根，且
  \[
  v(P_i)=\frac{1-3x(P_i)^2-2x(P_i)}{4x(P_i)}\qquad(i=2,3,4).
  \]
  \item 作为 \(E_0\) 上的亚纯微分，\(\dd f\) 的零极点除子满足
  \[
  \mathrm{div}(\dd f)=P_0+P_1+P_2+P_3+P_4-5O,
  \]
  从而在 \(E_0(\overline{\QQ})\) 的群律下有
  \[
  P_0+P_1+P_2+P_3+P_4=O,
  \qquad
  P_2+P_3+P_4=-(P_0+P_1)=\Bigl(\frac14,-\frac58\Bigr)\in E_0(\QQ).
  \]
\end{enumerate}
\end{proposition}
\begin{proof}
（1）在 \(E_0\) 的短 Weierstrass 配方 \(Y^2=X^3-X+\tfrac14\)（取 \(X=x\)、\(Y=v+\tfrac12\)）下，\(x\) 与 \(v\) 在 \(O\) 处的极点阶分别为 \(2\) 与 \(3\)，故 \(f=v+x^2\) 在 \(O\) 处极点阶为 \(4\)，从而 \((f)_\infty=4O\)。
\par
（2）对 \(E_0:v^2+v=x^3-x\) 微分得到
\[
(2v+1)\,\dd v=(3x^2-1)\,\dd x.
\]
于是
\[
\dd f=\dd(v+x^2)=\dd v+2x\,\dd x=\left(\frac{3x^2-1}{2v+1}+2x\right)\dd x.
\]
因此 \(\dd f=0\) 等价于 \(3x^2-1+2x(2v+1)=0\)，亦即
\[
4xv+3x^2+2x-1=0.
\]
在 \(x\neq 0\) 上可解得 \(v=(1-3x^2-2x)/(4x)\)；
将其代回 \(v^2+v=x^3-x\) 并化简即得 \((x^2-1)(16x^3-9x^2+1)=0\)，从而得到五个有限分歧点的显式刻画。
\par
（3）记 \(E_0\) 的标准正则微分
\[
\omega:=\frac{\dd x}{2v+1}.
\]
由上式可得 \(\dd f=(4xv+3x^2+2x-1)\,\omega\)。
由于 \(g(E_0)=1\)，\(\omega\) 在 \(E_0\) 上处处正则且无零极点；结合（2）知 \(\dd f\) 在五个有限分歧点处各有一个简单零点。
另一方面，由（1）可取局部参数 \(t\) 使 \(f\sim t^{-4}\)，则 \(\dd f\sim -4t^{-5}\dd t\)，故 \(\dd f\) 在 \(O\) 处为五阶极点，从而
\(\mathrm{div}(\dd f)=\sum_{i=0}^{4}P_i-5O\)。
该除子为典范类，而椭圆曲线的典范类在 \(\mathrm{Pic}^0(E_0)\) 中为零；在同一识别 \(E_0\simeq \mathrm{Pic}^0(E_0)\) 下得到群律恒等式 \(\sum_{i=0}^4 P_i=O\)。
最后，在短 Weierstrass 坐标 \((X,Y)\) 中，
\[
P_0=(1,-\tfrac12),\qquad P_1=(-1,\tfrac12),
\]
故弦割斜率为 \(-\tfrac12\)，从而 \(P_0+P_1=(\tfrac14,\tfrac18)\)，变回 \(v=Y-\tfrac12\) 得
\(
P_0+P_1=(\tfrac14,-\tfrac38)
\)；
取负元 \(v\mapsto -v-1\) 即得
\(
-(P_0+P_1)=(\tfrac14,-\tfrac58)
\)，从而 \(P_2+P_3+P_4=-(P_0+P_1)\)。
证毕。
\end{proof}
\par

\begin{theorem}[分歧型 \texorpdfstring{$(4;2^{5})$}{(4;2^5)} 覆盖的全对称单值化：\texorpdfstring{$\mathrm{Mon}(f)\cong S_4$}{Mon(f)=S4}]\label{thm:xi-terminal-zm-leyang-monodromy-s4}
沿用命题 \ref{prop:xi-terminal-zm-leyang-f-ramification} 的 \(f:E_0\to\PP^{1}\)。
则 \(f\) 的几何单值群同构于 \(S_4\)。
等价地，将 \(\Pi(\lambda,y)\) 视为 \(\QQ(y)\) 上关于 \(\lambda\) 的四次多项式，则其函数域伽罗瓦群为 \(S_4\)。
\end{theorem}
\begin{proof}
由命题 \ref{prop:xi-terminal-zm-leyang-f-ramification}，（i）表明绕 \(f=\infty\) 的局部单值为一个 \(4\)-循环；而任一有限分歧值处的局部单值为换位。
一个 \(4\)-循环与一个换位生成 \(S_4\)，故 \(\mathrm{Mon}(f)\cong S_4\)。
四次多项式的函数域伽罗瓦群与几何单值群之等价表述见结论 \ref{con:xi-terminal-zm-generic-galois-s4}。证毕。
\end{proof}
\par

\begin{corollary}[导子 \(37\) 的模性锚定与 \texorpdfstring{$\mathrm{GL}_2$}{GL2} 自守对象]\label{cor:xi-terminal-zm-leyang-langlands-level37}
令 \(\mathcal C\) 为定理 \ref{thm:xi-terminal-zm-leyang-elliptic-modularity-37} 的特征曲线（亦即 \(\Pi(\lambda,y)=0\) 的正规化），并令 \(E_0\) 为相应椭圆曲线。
则存在权 \(2\)、level \(37\) 的规范化新形式 \(g\in S_2(\Gamma_0(37))\) 及对应自守表示 \(\pi_{37}\)，使得
\[
L(\mathcal C,s)=L(E_0,s)=L(g,s)=L(\pi_{37},s).
\]
\end{corollary}
\begin{proof}
由定理 \ref{thm:xi-terminal-zm-leyang-elliptic-modularity-37}，\(\mathcal C\simeq E_0\) 且 \(E_0/\QQ\) 为椭圆曲线。
由模性定理存在权 \(2\) 新形式 \(g\) 使其 Hasse--Weil \(L\)-函数与 \(L(g,s)\) 一致。
对 level \(37\) 的外部命名与交叉核验接口见注记 \ref{rem:xi-terminal-elliptic-37a1-modular-anchor}。\cite{LMFDBEllipticCurve37a1,LMFDBNewform372aa,ConradTSWFinal}
证毕。
\end{proof}
\par

\begin{conjecture}[Kapustin--Witten 物理几何朗兰兹中的 Lee--Yang 分歧壁：范畴化对应]\label{conj:xi-terminal-zm-leyang-kapustin-witten-wall}
令 \(B\subset\PP^1\) 为定理 \ref{thm:xi-terminal-zm-leyang-monodromy-s4} 中覆盖 \(f:E_0\to\PP^1\) 的分歧值集合，并令 \(U:=\PP^1\setminus B\)。
记
\[
\mathcal L:=\left.(f_\ast\underline{\QQ})\right|_{U}
\]
为由四次覆盖推前得到的秩 \(4\) 局部系统，则其几何单值像为 \(S_4\)。
猜想在 Kapustin--Witten 所述的 \(4\mathrm d\ \mathcal N=4\) 超对称 Yang--Mills 的拓扑扭曲与边界/线算符框架下，可构造一类范畴对象 \(\mathfrak E(\mathcal L)\)，使其在几何朗兰兹等价中成为 Hecke 特征对象，特征值局部系统为 \(\mathcal L\)，并且 \(B\) 中的有限分歧值对应于耦合常数复延拓中的壁穿越与非正则极限。\cite{KapustinWitten2006GeometricLanglands}
\end{conjecture}

